% 原书物理页=印刷页55--59；2.3 全节.（2.3.4 段命题证明的后半部分、其后的注与 Cartan 判据推论实际印于原书第 60 页上部、2.4 节标题之前，属本节内容，一并译出.）
\section{幂零与可解 Lie 代数}\label{sec:lie-2-3}

设 $\mathfrak g$ 是一个 $K$-Lie 代数. 对 $\mathfrak g$ 的两个理想 $\mathfrak h_{1},\mathfrak h_{2}$，我们定义
\[
[\mathfrak h_{1},\mathfrak h_{2}]=\{\text{由 }[x_{1},x_{2}]\ (x_{1}\in\mathfrak h_{1},\ x_{2}\in\mathfrak h_{2})\ \text{生成的 }K\text{-向量空间}\}.
\]
根据 Jacobi 恒等式，它仍是 $\mathfrak g$ 的一个理想.
\begin{example*}
（\newterm{lie-derived-ideal}{导理想}{derived ideal}）$\mathfrak g$ 的理想 $[\mathfrak g,\mathfrak g]$ 称为 $\mathfrak g$ 的导理想. 它为零当且仅当 $\mathfrak g$ 是阿贝尔的. 一般地，它含于从 $\mathfrak g$ 到任何阿贝尔 Lie 代数的每个态射之中\footnote{译注：原文如此；按上下文意指它含于每个此类态射的核之中.}，因此 $\mathfrak g_{\mathrm{ab}}:=\mathfrak g/[\mathfrak g,\mathfrak g]$ 是 $\mathfrak g$ 的\emph{最大阿贝尔商}.
\end{example*}
更一般地，我们定义理想的如下两个（按包含关系）递降序列：
\begin{enumerate}[label=\roman*)]
\item 由 $C^{0}\mathfrak g:=\mathfrak g$ 与 $C^{n+1}\mathfrak g:=[\mathfrak g,C^{n}\mathfrak g]$（对 $n\geqslant 0$）定义的\newterm{lie-lower-central-series}{降中心列}{lower central series}；
\item 由 $D^{0}\mathfrak g:=\mathfrak g$ 与 $D^{n+1}\mathfrak g:=[D^{n}\mathfrak g,D^{n}\mathfrak g]$（对 $n\geqslant 0$）定义的\newterm{lie-derived-series}{导列}{derived series}.
\end{enumerate}
我们有 $C^{1}\mathfrak g=D^{1}\mathfrak g=[\mathfrak g,\mathfrak g]$，并且当 $n\geqslant 1$ 时有 $C^{n}\mathfrak g\supset D^{n}\mathfrak g$.
\begin{definition}\label{def:lie-2-3-1}
称 $\mathfrak g$ 是
\begin{itemize}[label=--]
\item \newterm{lie-nilpotent}{幂零}{nilpotent}的，如果存在 $n\in\mathbb N$ 使得 $C^{n}\mathfrak g=0$；
\item \newterm{lie-solvable}{可解}{solvable}的，如果存在 $n\in\mathbb N$ 使得 $D^{n}\mathfrak g=0$.
\end{itemize}
\end{definition}
于是我们有``阿贝尔''$\Rightarrow$``幂零''$\Rightarrow$``可解''. 读者还会注意到，这里与有限群的幂零性和可解性概念之间存在平行关系.
\begin{exercise*}
设 $\mathfrak g$ 是一个 $K$-Lie 代数.
\begin{enumerate}[label=\roman*)]
\item
\begin{enumerate}[label=\alph*)]
\item 证明 $\mathfrak g$ 是幂零的当且仅当存在 $\mathfrak g$ 的一个理想序列 $\mathfrak g=\mathfrak h_{0}\supset\mathfrak h_{2}\supset\ldots\supset\mathfrak h_{n}=0$\footnote{译注：原书如此，按上下文应为 $\mathfrak h_{1}$.}，使得对一切 $i$ 有 $[\mathfrak g,\mathfrak h_{i}]\subset\mathfrak h_{i+1}$.
\item 由此推出，幂零 Lie 代数的任何 Lie 子代数与任何商代数都是幂零的.
\item 证明 $\mathfrak g$ 是幂零的当且仅当 $\operatorname{ad}(\mathfrak g)=\mathfrak g/\mathfrak z(\mathfrak g)$ 是幂零的.
\end{enumerate}
\item
\begin{enumerate}[label=\alph*)]
\item 证明 $\mathfrak g$ 是可解的当且仅当存在 $\mathfrak g$ 的一个理想序列 $\mathfrak g=\mathfrak h_{0}\supset\mathfrak h_{2}\supset\ldots\supset\mathfrak h_{n}=0$\footnote{译注：原书如此，按上下文应为 $\mathfrak h_{1}$.}，使得对一切 $i$，$\mathfrak h_{i}/\mathfrak h_{i+1}$ 都是阿贝尔的.
\item 由此推出，若 $\mathfrak h$ 是 $\mathfrak g$ 的一个理想，则 $\mathfrak g$ 是可解的当且仅当 $\mathfrak h$ 与 $\mathfrak g/\mathfrak h$ 都是可解的.\footnote{注意，``若''这一方向对``幂零''并不成立.}
\end{enumerate}
\end{enumerate}
\end{exercise*}
\begin{definition}\label{def:lie-2-3-2}
设 $V$ 是一个维数为 $n$ 的 $K$-向量空间. $V$ 的一个\newterm{lie-complete-flag}{完全旗}{complete flag}是一个由子空间组成的递增序列 $\mathcal V=(0=\mathcal V_{0}\subset\mathcal V_{1}\subset\cdots\subset\mathcal V_{n}=V)$，其中 $\dim(\mathcal V_{i})=i$. 我们给它配上 $\mathfrak{gl}(V)$ 的两个 Lie 子代数：
\begin{gather*}
\mathfrak b(\mathcal V):=\{X\in\mathfrak{gl}(V),\ X(\mathcal V_{i})\subset\mathcal V_{i},\ \text{对 }i=1,\cdots,n.\}\quad\text{以及}\\
\mathfrak b^{\mathrm{nilp}}(\mathcal V):=\{X\in\mathfrak{gl}(V),\ X(\mathcal V_{i})\subset\mathcal V_{i-1},\ \text{对 }i=1,\cdots,n.\}.
\end{gather*}
\end{definition}
\begin{remark*}
若选取一组基 $v_{1},\cdots,v_{n}$ 使得 $\mathcal V_{i}=\operatorname{vect}(v_{1},\cdots,v_{i})$，则由此得到的同构 $\mathfrak{gl}(V)\xrightarrow{\ \sim\ }\mathfrak{gl}_{n}(K)$ 把 $\mathfrak b(\mathcal V)$、相应地把 $\mathfrak b^{\mathrm{nilp}}(\mathcal V)$ 等同于 $\mathrm M_{n}(K)$ 中的上三角矩阵、相应地严格上三角矩阵所构成的空间.
\end{remark*}
\begin{exercise*}
验证 $\mathfrak b(\mathcal V)$ 确实是 $\mathfrak{gl}(V)$ 的一个 Lie 子代数，然后证明：
\begin{enumerate}[label=\roman*)]
\item $[\mathfrak b(\mathcal V),\mathfrak b(\mathcal V)]=\mathfrak b^{\mathrm{nilp}}(\mathcal V)$，且 $[\mathfrak b(\mathcal V),\mathfrak b^{\mathrm{nilp}}(\mathcal V)]=\mathfrak b^{\mathrm{nilp}}(\mathcal V)$.
\item 对 $k\geqslant 1$，$C^{k}(\mathfrak b^{\mathrm{nilp}}(\mathcal V))=\{X\in\mathfrak{gl}(V),\ X(\mathcal V_{i})\subset\mathcal V_{i-k},\ \text{对 }i=1,\cdots,n.\}$
\end{enumerate}
由此推出 $\mathfrak b^{\mathrm{nilp}}(\mathcal V)$ 是幂零的，而 $\mathfrak b(\mathcal V)$ 是可解的但不是幂零的.
\end{exercise*}
\begin{numpara}[表示论刻画]\label{np:lie-2-3-3}
本段中我们假定 $K$ 是\newterm{lie-algebraically-closed}{代数闭}{algebraically closed}的且特征为零，而且局限于 $K=\mathbb C$ 并无妨. 我们首先要看到，上面的可解代数的例子是典型的.
\begin{theorem*}[Lie]
设 $\mathfrak g$ 是一个 $K$-Lie 代数，其中 $K$ 代数闭. 则下列诸条等价：
\begin{enumerate}[label=\roman*)]
\item $\mathfrak g$ 是可解的.
\item $\mathfrak g$ 的任何不可约表示都是一维的.
\item[$\mathrm{ii})'$] $\mathfrak g$ 的任何表示都含有对 $\mathfrak g$ 公共的特征向量.
\item $\mathfrak g$ 的任何表示 $(V,r)$ 都是\newterm{lie-trigonalizable}{可三角化}{trigonalizable}的，即存在 $V$ 的一个\newterm{lie-flag}{旗}{flag} $\mathcal V$ 使得 $r(\mathfrak g)\subset\mathfrak b(\mathcal V)$.
\end{enumerate}
\end{theorem*}
\begin{proof}
等价 ii)$\Leftrightarrow$ii)$'$ 是显然的，蕴涵 iii)$\Rightarrow$ii)$'$ 亦然，而蕴涵 ii)$'\Rightarrow$iii) 可由一个显然的归纳看出.
现证 iii)$\Rightarrow$i). 把 iii) 应用于伴随表示. 我们得到一个由子空间构成的完全旗 $0\subset\mathfrak h_{1}\subset\cdots\subset\mathfrak h_{n}=\mathfrak g$，使得 $\operatorname{ad}_{\mathfrak g}(\mathfrak h_{i})\subset\mathfrak h_{i}$，这表明诸 $\mathfrak h_{i}$ 都是 $\mathfrak g$ 的理想. 由于商 $\mathfrak h_{i}/\mathfrak h_{i+1}$ 的维数为 $1$，它是一个阿贝尔 Lie 代数，于是可以应用上面习题中的判据，从而推出 $\mathfrak g$ 是可解的.
现在证 i)$\Rightarrow$ii)$'$，对 $\dim(\mathfrak g)$ 作归纳. 若 $\dim(\mathfrak g)=1$，结果是显然的，因为有限维 $K$-向量空间的任何自同态都具有特征向量（$K$ 代数闭）. 设 $\dim(\mathfrak g)>1$，并设 $\mathfrak h$ 是 $\mathfrak g$ 中包含 $D\mathfrak g$ 的一个超平面（注意 $\mathfrak g/D\mathfrak g\neq 0$）. 这时 $\mathfrak h$ 是 $\mathfrak g$ 的理想，并且由归纳假设，存在 $V$ 的一个对 $\mathfrak h$ 的所有元素公共的特征向量 $v$. 设 $\lambda$ 是 $\mathfrak h$ 上由``对 $\mathfrak h$ 中一切 $X$ 有 $X.v=\lambda(X)v$''定义的线性形式. 令
\[
V_{\lambda}:=\{v\in V,\ \forall X\in\mathfrak h,\ X.v=\lambda(X)v\},
\]
即 $V$ 在 $\mathfrak h$ 作用下的``$\lambda$-特征''子空间，它因而非零. 我们暂且承认 $V_{\lambda}$ 在 $\mathfrak g$ 下稳定. 设 $Y$ 是 $\mathfrak g$ 中 $\mathfrak h$ 的一个补直线的生成元. 由于 $K$ 代数闭，$Y$ 具有一个属于 $V_{\lambda}$ 的特征向量 $v'$. 于是，$v'$ 既对 $\mathfrak h$ 是特征向量，对整个 $\mathfrak g$ 也是特征向量，定理得证.
尚需证明 $V_{\lambda}$ 在 $\mathfrak g$ 下稳定. 由等式
\[
XYv=YXv+[X,Y]v=\lambda(X)Yv+\lambda([X,Y])v
\]
对一切 $X\in\mathfrak h$、$Y\in\mathfrak g$、$v\in V_{\lambda}$ 成立，我们把问题化归为证明：对一切 $X\in\mathfrak h$、$Y\in\mathfrak g$ 有 $\lambda([X,Y])=0$. 为此，取定 $V_{\lambda}$ 中一个向量 $w\neq 0$，并记
\[
W_{k}:=\operatorname{vect}\{w,\,Yw,\,\cdots,\,Y^{k}w\}
\]
公式 $XY^{k}w=YXY^{(k-1)}w+[X,Y]Y^{(k-1)}w$ 使我们能够用归纳法证明
\[
\forall X\in\mathfrak h,\ XW_{k}\subset W_{k}\quad\text{且}\quad(\ast):\ XY^{k}w\in\lambda(X)Y^{k}w+W_{k-1}.
\]
然后考虑诸 $W_{k}$ 的并 $W$. 它是一个在 $\mathfrak h$ 下稳定的子空间，且 $(\ast)$ 表明对一切 $X\in\mathfrak h$ 有 $\operatorname{tr}(X|_{W})=\dim(W)\,\lambda(X)$. 特别地，$\dim(W)\,\lambda([X,Y])=\operatorname{tr}([X,Y]|_{W})$. 而由构造，$W$ 也在 $Y$ 下稳定，于是 $\operatorname{tr}([X,Y]|_{W})=\operatorname{tr}([X|_{W},Y|_{W}])=0$. 因此如所愿地得到 $\lambda([X,Y])=0$（因为 $\dim(W)$ 在特征为零的域 $K$ 中非零）.
\end{proof}
\begin{remark*}
该定理推广了线性代数中的经典结果：两两交换的一族自同态可同时三角化.
\end{remark*}
为了通过表示来刻画幂零代数，我们引入几个定义与记号. 若 $\lambda:\mathfrak g\longrightarrow K$ 是一个映射，$(V,r)$ 是一个表示，则记
\[
V_{\lambda}:=\{v\in V,\ \forall X\in\mathfrak g,\ r(X)v=\lambda(X)v\}
\]
为相应的特征空间，而
\[
V^{\lambda}:=\{v\in V,\ \exists n\in\mathbb N,\ \forall X\in\mathfrak g,\ (r(X)-\lambda(X))^{n}v=0\}
\]
为相应的``广义''特征空间.
\begin{definition*}
若 $V_{\lambda}\neq 0$，则称 $\lambda$ 是 $\mathfrak g$ 在表示 $(V,r)$ 中的一个权. 此时 $V_{\lambda}$ 称为权 $\lambda$ 的\newterm{lie-weight-space}{权空间}{weight space}，$V^{\lambda}$ 称为广义权空间.
\end{definition*}
\begin{exercise*}
若 $\lambda$ 是某个表示的一个权，则 $\lambda$ 是 $K$-线性的，且 $\lambda([\mathfrak g,\mathfrak g])=0$. 因此它是 $\mathfrak g$ 的\newterm{lie-abelianization}{阿贝尔化}{abelianization} $\mathfrak g_{\mathrm{ab}}$ 的对偶 $\mathfrak g_{\mathrm{ab}}^{*}$ 中的一个元素.
再注意，一般说来并不清楚 $V^{\lambda}$ 是否在 $\mathfrak g$ 下稳定.
\end{exercise*}
\begin{proposition*}[Engel 定理]
若 $V=V^{\lambda}$，则 $\lambda\in\mathfrak g_{\mathrm{ab}}^{*}$，且存在 $V$ 的一个完全旗 $\mathcal V$ 使得 $(r-\lambda)(\mathfrak g)\subset\mathfrak b^{\mathrm{nilp}}(\mathcal V)$.
\end{proposition*}
\begin{remark*}
具体地说，这表明：若对每个 $X$ 都存在一组基，使 $r(X)$ 在这组基下是上三角的且对角线上为诸 $\lambda(X)$，则事实上存在一组基，使\emph{所有} $r(X)$ 都具有这种形状.
\end{remark*}
\begin{proof}
由假设，对一切 $X\in\mathfrak g$ 有 $\operatorname{tr}(r(X))=\lambda(X)\dim(V)$. 这表明（需要 $K$ 的特征为 $0$，以保证 $\dim(V)\neq 0$ 于 $K$ 中）$\lambda$ 是 $K$-线性的且在 $[\mathfrak g,\mathfrak g]$ 上为零，故来自 $\mathfrak g_{\mathrm{ab}}^{*}$. 但这样一来，映射 $X\mapsto r(X)-\lambda(X)\,\mathrm{id}_{V}$ 便在同一个空间 $V$ 上定义了 $\mathfrak g$ 的一个表示. 不妨用 $r-\lambda$ 替换 $r$，从而可以设 $\lambda=0$. 通过一个归纳论证，我们只需证明：
\[
(\ast)\quad\text{对每个使 }r(X)\ \text{对一切 }X\ \text{幂零的表示 }(V,r)\text{，有 }V^{\mathfrak g}\neq 0.
\]
我们将对 $\dim(\mathfrak g)$ 作归纳来完成它. 若 $\dim(\mathfrak g)=1$，则对 $\mathfrak g$ 中每个非零的 $X$ 有 $V^{\mathfrak g}=\operatorname{Ker}(r(X))$. 由 $r(X)$ 幂零的假设知其核非零，正如所愿. 现设 $\dim(\mathfrak g)>1$，并设我们已有一个非平凡的理想 $\mathfrak h$，即它异于 $\{0\}$ 与 $\mathfrak g$. 这时，应用于 $\mathfrak h$ 的归纳假设告诉我们 $V^{\mathfrak h}\neq\{0\}$. 此外，$V^{\mathfrak h}$ 在 $\mathfrak g$ 的作用下稳定，而该作用经由 $\mathfrak g/\mathfrak h$ 分解. 于是，应用于 $\mathfrak g/\mathfrak h$ 的归纳假设告诉我们 $(V^{\mathfrak h})^{\mathfrak g/\mathfrak h}\neq 0$. 然而 $V^{\mathfrak g}=(V^{\mathfrak h})^{\mathfrak g/\mathfrak h}$.
故尚需证明非平凡理想 $\mathfrak h$ 的存在性. 不妨用 $r(\mathfrak g)$ 替换 $\mathfrak g$，从而可设 $\mathfrak g\subset\mathfrak{gl}(V)$. 这时选取 $\mathfrak g$ 的一个维数最大的真 Lie 子代数 $\mathfrak h$. $\mathfrak g$ 的伴随表示 $\operatorname{ad}$ 限制到 $\mathfrak h$ 上诱导出 $\mathfrak h$ 在商 $\mathfrak g/\mathfrak h$ 上的一个表示. $\mathfrak h$ 在该表示下的像由幂零自同态组成. 事实上，根据下面的习题，对一切 $X\in\mathfrak g$，$\mathfrak g$ 上的自同态 $\operatorname{ad}(X)$ 是幂零的. 于是可以应用我们的归纳假设推出：在 $\mathfrak g/\mathfrak h$ 中存在被 $\mathfrak h$ 的作用零化的元素 $\overline{X}\neq 0$. 设 $X\in\mathfrak g$ 是位于 $\overline{X}$ 之上的一个元素. 于是 $X\notin\mathfrak h$，且 $[\mathfrak h,X]=\operatorname{ad}(\mathfrak h)(X)\subseteq\mathfrak h$. 这蕴含 $\mathfrak h\oplus K.X$ 是一个以 $\mathfrak h$ 为理想的 Lie 代数. 而由 $\mathfrak h$ 的极大性，$\mathfrak h\oplus K.X=\mathfrak g$. 这样，$\mathfrak h$ 是 $\mathfrak g$ 的余维数为 $1$ 的理想.
\end{proof}
\begin{exercise*}
设 $A$ 是一个 $K$-结合代数，并设 $\operatorname{ad}(x)(y)=[x,y]:=xy-yx$（$x,y\in A$）. 证明 $\operatorname{ad}(x)^{n}(y)\in\operatorname{Vect}_{K}\{x^{k}yx^{n-k},\ k=0,\cdots,n\}$.
\end{exercise*}
\begin{corollary*}
设 $\mathfrak g$ 是一个有限维 $K$-Lie 代数. 则 $\mathfrak g$ 是幂零的当且仅当对一切 $X\in\mathfrak g$，$\operatorname{ad}_{X}$ 是幂零的.
\end{corollary*}
\begin{proof}
由于 $\operatorname{ad}_{X}^{n}(\mathfrak g)\subset C^{n}\mathfrak g$，可见该条件是必要的. 反之，设对一切 $X\in\mathfrak g$，$\operatorname{ad}_{X}$ 都是幂零的. 根据上面的命题（取 $\lambda=0$），存在一个旗 $\mathfrak h_{1}\subset\mathfrak h_{2}\subset\cdots\,\mathfrak h_{n}=\mathfrak h$\footnote{译注：原书如此，按上下文应为 $\mathfrak h_{n}=\mathfrak g$.}，使得对一切 $X$ 有 $\operatorname{ad}_{X}(\mathfrak h_{i})\subset\mathfrak h_{i-1}$，亦即 $[\mathfrak g,\mathfrak h_{i}]\subset\mathfrak h_{i-1}$. 把诸 $\mathfrak h_{i}$ 重编号为 $\mathfrak h_{n-i}$ 之后，我们便认出这就是定义之后的那条幂零性判据.
\end{proof}
\begin{lemma*}
设 $\mathfrak g$ 是一个幂零 $K$-Lie 代数，$\mathfrak h\subset\mathfrak g$ 是一个 Lie 子代数，$\lambda:\mathfrak h\longrightarrow K$ 是一个映射，$V^{\lambda}$ 是广义特征空间. 则 $V^{\lambda}$ 在 $\mathfrak g$ 下稳定.
\end{lemma*}
\begin{proof}
设 $Y\in\mathfrak g$，$X\in\mathfrak h$，$v\in V^{\lambda}$. 我们要证明
\[
(r(X)-\lambda(X))^{n}r(Y)v=0\quad\text{对足够大的 }n.
\]
而对于 $V$ 的两个自同态 $f,g$，我们可以用归纳法（习题）证明\footnote{译注：原书此式中的二项式系数误印作 $\tbinom{k}{n}$（上下颠倒；对照 2.5.1 中同一公式），应为 $\tbinom{n}{k}$，此处按正确记法译出.}
\[
f^{n}\circ g=\sum_{k=0}^{n}\binom{n}{k}(\operatorname{ad}_{f})^{k}(g)\circ f^{n-k}.
\]
令 $f=r(X)-\lambda(X)$，$g=r(Y)$. 选取 $m$ 使得 $f^{m}(v)=0$ 且 $\operatorname{ad}_{f}^{m}(g)=0$（事实上 $\operatorname{ad}_{f}=\operatorname{ad}_{r(X)}$，而 $\mathfrak g$ 幂零，故 $r(\mathfrak g)$ 幂零）ootnote{译注：原文括号内印作``ad$_f=$ad$_{r(X)}$ et $\mathfrak g$, donc $r(\mathfrak g)$ est nilpotente''，其中``et $\mathfrak g$''语义未完；按上下文补足为``$\mathfrak g$ 幂零''.}. 那么上面的公式便给出 $f^{2m}\circ g(v)=0$. 这样我们就证明了 $V^{\lambda}$ 在 $\mathfrak g$ 下稳定.
\end{proof}
\begin{theorem*}
设 $K$ 代数闭，并设 $\mathfrak g$ 是一个有限维 $K$-Lie 代数. 则下列诸条等价：
\begin{enumerate}[label=\roman*)]
\item $\mathfrak g$ 是幂零的；
\item 任何表示 $(V,r)$ 都分解为 $V=\bigoplus_{\lambda}V^{\lambda}$，其中 $\lambda$ 取遍 $\mathfrak g$ 在 $V$ 中的权，且每个 $V^{\lambda}$ 都在 $\mathfrak g$ 下稳定.
\end{enumerate}
\end{theorem*}
\begin{proof}
ii)$\Rightarrow$i). 把 ii) 应用于伴随表示. 我们由此得到 $\mathfrak g=\bigoplus_{\lambda}\mathfrak g^{\lambda}$ 这一理想的直和分解. 证明 $\mathfrak g^{\lambda}\neq 0\Rightarrow\lambda=0$. 事实上，设 $Y\in\mathfrak g^{\lambda}$. 自同态 $(\operatorname{ad}_{Y}-\lambda(Y))|_{\mathfrak g^{\lambda}}$ 是幂零的，而 $(\operatorname{ad}_{Y})|_{\mathfrak g^{\lambda}}$ 不可逆，因为 $\operatorname{ad}_{Y}(Y)=0$. 由此推出 $\lambda(Y)=0$，从而 $\lambda|_{\mathfrak g^{\lambda}}=0$. 但由于 $[\mathfrak g^{\lambda'},\mathfrak g^{\lambda}]\subset\mathfrak g^{\lambda'}\cap\mathfrak g^{\lambda}$，对其余的权也有 $\lambda|_{\mathfrak g^{\lambda'}}=0$，故最终 $\lambda=0$. 根据上面的推论，$\mathfrak g$ 是幂零的.
i)$\Rightarrow$ii). 对 $\dim(V)$ 作归纳. 若所有的 $r(X)$ 都只有一个特征值，则对某个函数 $\lambda$ 有 $V=V^{\lambda}$，而且我们已经看到 $\lambda$ 是一个权（上面的命题）. 否则，存在一个至少具有两个不同特征值的 $X$（$K$ 代数闭）. 由此得到 $r(X)$ 的广义特征空间的一个非平凡分解 $V=V_{X}^{\lambda_{1}}\oplus\cdots\oplus V_{X}^{\lambda_{k}}$. 根据上面的引理，每个 $V_{X}^{\lambda_{i}}$ 都在 $\mathfrak g$ 下稳定，于是可以对它应用归纳假设.
\end{proof}
\begin{remark*}
该定理对 $r(\mathfrak g)$ 的元素给出了某种``同时 Jordan 形式''.
\end{remark*}
\end{numpara}
\begin{numpara}[用 Killing 型刻画]\label{np:lie-2-3-4}
仍设 $K$ 代数闭且特征为零.
\begin{proposition*}
设 $\mathfrak g$ 是 $\mathfrak{gl}_{n}(K)$ 的一个 $K$-Lie 子代数，满足对一切 $X,Y\in\mathfrak g$ 有 $\operatorname{tr}(XY)=0$. 那么 $[\mathfrak g,\mathfrak g]$ 是幂零的，从而 $\mathfrak g$ 是可解的.
\end{proposition*}
\begin{remark*}
$\mathfrak g$ 的可解性由定义~\ref{def:lie-2-3-1} 之后的习题之 ii)(b) 应用于 $\mathfrak h=[\mathfrak g,\mathfrak g]$ 与 $\mathfrak g/\mathfrak h=\mathfrak g_{\mathrm{ab}}$（它是阿贝尔的）而得到.
\end{remark*}
\begin{proof}
为简单起见，我们将假定 $K=\mathbb C$ 或 $K=\overline{\mathbb Q}\subset\mathbb C$（注意，若 $K$ 的基数至多为连续统的基数（at most second-countable），则众所周知（在承认不可数选择公理的前提下）$K$ 实际上同构于 $\mathbb C$ 或 $\overline{\mathbb Q}\subset\mathbb C$）. 根据 Engel 定理，只需证明 $[\mathfrak g,\mathfrak g]\subset\mathfrak{gl}_{n}(K)$ 中的任何元素 $X$ 都是幂零的，也就是说，它的特征值 $\lambda_{1},\cdots,\lambda_{n}$ 全为零. 选取一个多项式 $f\in K[T]$ 使得对一切 $i$ 有 $f(\lambda_{i})=\overline{\lambda_{i}}$（例如一个 Lagrange 插值多项式）. 设 $X=X_{s}+X_{n}$ 是 $X$ 的 Jordan 分解. 那么 $\operatorname{tr}(f(X_{s})X)=\sum_{i}\lambda_{i}\overline{\lambda_{i}}=\sum_{i}|\lambda_{i}|^{2}$. 因此我们只需证明 $\operatorname{tr}(f(X_{s})X)=0$. 由于 $X\in[\mathfrak g,\mathfrak g]$ 是若干 $[Y,Z]$（$Y,Z\in\mathfrak g$）的线性组合，并且
\[
\operatorname{tr}(f(X_{s})[Y,Z])=\operatorname{tr}([f(X_{s}),Y]Z),
\]
故只需证明 $\operatorname{ad}_{f(X_{s})}(\mathfrak g)\subset\mathfrak g$，以便能够应用假设. 然而，$\mathfrak{gl}_{n}(K)$ 的自同态 $\operatorname{ad}_{f(X_{s})}$ 与自同态 $\operatorname{ad}_{X_{s}}$ 在同一组基下可对角化. 更确切地说，若 $(e_{1},\cdots,e_{n})$ 是 $K^{n}$ 的一组使 $X_{s}$ 对角化的基（特征值为 $\lambda_{1},\cdots,\lambda_{n}$），而 $(E_{ij})_{1\leqslant i,j\leqslant n}$ 是 $\mathfrak{gl}_{n}(K)$ 的相应基，则 $E_{ij}$ 是 $\operatorname{ad}_{X_{s}}$、相应地 $\operatorname{ad}_{f(X_{s})}$ 的属于特征值 $\lambda_{i}-\lambda_{j}$、相应地 $f(\lambda_{i})-f(\lambda_{j})=\lambda_{i}-\lambda_{j}$ 的特征向量. 由此可见，对每个满足 $g(\lambda_{i}-\lambda_{j})=\lambda_{i}-\lambda_{j}$ 的多项式 $g\in K[T]$（这样的多项式是存在的，例如一个插值多项式），有 $\operatorname{ad}_{f(X_{s})}=g(\operatorname{ad}_{X_{s}})$. 因此我们只需证明 $\operatorname{ad}_{X_{s}}(\mathfrak g)\subset\mathfrak g$. 然而，$\operatorname{ad}_{X_{s}}$ 是 $\operatorname{ad}_{X}$ 在其 Jordan 分解中的半单部分（见下面的注），从而它是 $\operatorname{ad}_{X}$ 的一个多项式，并且由于 $\operatorname{ad}_{X}(\mathfrak g)\subset\mathfrak g$，证明完毕.
\end{proof}
\begin{remark*}
我们用到了这样一个事实：若 $X=X_{s}+X_{n}$ 是 $K$-向量空间 $V$ 上一个 $K$-线性自同态的 Jordan 分解，则 $\operatorname{ad}_{X}=\operatorname{ad}_{X_{s}}+\operatorname{ad}_{X_{n}}$ 是 $\mathfrak{gl}(V)$ 的自同态 $\operatorname{ad}_{X}$ 的 Jordan 分解. 事实上，我们已经注意到 $\operatorname{ad}_{X_{n}}$ 是幂零的（Engel 定理后面的习题），而我们在上面的证明中刚刚解释了 $\operatorname{ad}_{X_{s}}$ 是可对角化的. 因此只需注意 $\operatorname{ad}_{X_{s}}$ 与 $\operatorname{ad}_{X_{n}}$ 交换.
\end{remark*}
\begin{corollary*}[Cartan 判据]
设 $\mathfrak g$ 是一个有限维 $K$-Lie 代数\footnote{译注：原文作 ``une $K$-algèbre de dimension finie''，未写明 Lie，按上下文应为 Lie 代数.}. 则 $\mathfrak g$ 是可解的当且仅当 $B_{\operatorname{ad}}(\mathfrak g,[\mathfrak g,\mathfrak g])=0$.
\end{corollary*}
\begin{proof}
$\Rightarrow$. 若 $\mathfrak g$ 可解，Lie 定理告诉我们存在 $\mathfrak g$ 中的一个旗 $\mathcal V=(\mathfrak h_{1}\subset\cdots\subset\mathfrak h_{n})$ 使得 $\operatorname{ad}(\mathfrak g)\subset\mathfrak b(\mathcal V)$. 于是 $\operatorname{ad}([\mathfrak g,\mathfrak g])\subset[\operatorname{ad}(\mathfrak g),\operatorname{ad}(\mathfrak g)]\subset\mathfrak b^{\mathrm{nilp}}(\mathcal V)$，从而对一切 $X\in\mathfrak g$ 与 $Y\in[\mathfrak g,\mathfrak g]$ 有 $\operatorname{ad}_{X}\operatorname{ad}_{Y}\in\mathfrak b^{\mathrm{nilp}}(\mathcal V)$，因此 $\operatorname{tr}(\operatorname{ad}_{X}\operatorname{ad}_{Y})=0$.
$\Leftarrow$. 设 $\mathfrak h:=\operatorname{ad}([\mathfrak g,\mathfrak g])$，其中 $\operatorname{ad}:\mathfrak g\longrightarrow\mathfrak{gl}(\mathfrak g)$ 是 $\mathfrak g$ 的伴随表示. 由于 $\mathfrak h\simeq[\mathfrak g,\mathfrak g]/(\mathfrak z(\mathfrak g)\cap[\mathfrak g,\mathfrak g])$ 且 $\mathfrak g/[\mathfrak g,\mathfrak g]$ 是阿贝尔的，可见 $\mathfrak g$ 可解当且仅当 $\mathfrak h$ 可解（见定义~\ref{def:lie-2-3-1} 之后的习题之 ii)(b)）. 而 $\mathfrak h$ 的可解性由前面的命题得到.
\end{proof}
\end{numpara}
